\chapter{A-infinity Categories}

\begin{definition}[\Ainf{}-precategory]
\lean{AInfinityCategoryTheory.AInfinityPreCategory}
\label{AInfinityCategoryTheory.AInfinityPreCategory}
\leanok
\uses{AInfinityTheory.RLinearGQuiver, AInfinityTheory.composableHomType, AInfinityTheory.operationTargetType}
An \Ainf{}-precategory over $R$ with object type $\mathrm{Obj}$ consists of a
graded $R$-linear quiver together with operations
\[
  m_n
\]
of the prescribed multilinear type for every positive arity $n$.
\end{definition}

\section{Chains}

\begin{definition}[Chain]
\lean{AInfinityCategoryTheory.AInfinityPreCategory.Chain}
\label{AInfinityCategoryTheory.AInfinityPreCategory.Chain}
\leanok
\uses{AInfinityCategoryTheory.AInfinityPreCategory}
A chain in an \Ainf{}-precategory consists of a positive length $n$, a string of
$n+1$ objects, and a degree assigned to each of the $n$ composable morphism
slots.
\end{definition}

\begin{definition}[Chain source]
\lean{AInfinityCategoryTheory.AInfinityPreCategory.Chain.source}
\label{AInfinityCategoryTheory.AInfinityPreCategory.Chain.source}
\leanok
\uses{AInfinityCategoryTheory.AInfinityPreCategory.Chain}
The source of a chain is its initial object.
\end{definition}

\begin{definition}[Chain target]
\lean{AInfinityCategoryTheory.AInfinityPreCategory.Chain.target}
\label{AInfinityCategoryTheory.AInfinityPreCategory.Chain.target}
\leanok
\uses{AInfinityCategoryTheory.AInfinityPreCategory.Chain}
The target of a chain is its final object.
\end{definition}

\begin{definition}[Chain link type]
\lean{AInfinityCategoryTheory.AInfinityPreCategory.Chain.link}
\label{AInfinityCategoryTheory.AInfinityPreCategory.Chain.link}
\leanok
\uses{AInfinityCategoryTheory.AInfinityPreCategory.Chain, AInfinityTheory.composableHomType}
For a chain and an index $i$,
\lean{AInfinityCategoryTheory.AInfinityPreCategory.Chain.link} is the module of
the $i$-th composable morphism in that chain.
\end{definition}

\begin{definition}[Chain operation target]
\lean{AInfinityCategoryTheory.AInfinityPreCategory.Chain.operationTarget}
\label{AInfinityCategoryTheory.AInfinityPreCategory.Chain.operationTarget}
\leanok
\uses{AInfinityCategoryTheory.AInfinityPreCategory.Chain, AInfinityTheory.operationTargetType}
For a chain, this is the graded hom-module that receives the corresponding
multilinear operation.
\end{definition}

\section{Stasheff Identities}

\begin{definition}[Precategory Stasheff sum]
\lean{AInfinityCategoryTheory.AInfinityPreCategory.stasheffSum}
\label{AInfinityCategoryTheory.AInfinityPreCategory.stasheffSum}
\leanok
\uses{AInfinityCategoryTheory.AInfinityPreCategory, AInfinityTheory.indexedStasheffSum}
The Stasheff sum of an \Ainf{}-precategory is the indexed Stasheff sum formed
from its structure maps.
\end{definition}

\begin{definition}[Precategory Stasheff property]
\lean{AInfinityCategoryTheory.AInfinityPreCategory.satisfiesStasheff}
\label{AInfinityCategoryTheory.AInfinityPreCategory.satisfiesStasheff}
\leanok
\uses{AInfinityCategoryTheory.AInfinityPreCategory, AInfinityTheory.indexedSatisfiesStasheff}
An \Ainf{}-precategory satisfies the Stasheff identities when its structure maps
define an object-indexed system satisfying
\lean{AInfinityTheory.indexedSatisfiesStasheff}.
\end{definition}

\begin{definition}[\Ainf{}-category]
\lean{AInfinityCategoryTheory.AInfinityCategory}
\label{AInfinityCategoryTheory.AInfinityCategory}
\leanok
\uses{AInfinityCategoryTheory.AInfinityPreCategory, AInfinityCategoryTheory.AInfinityPreCategory.satisfiesStasheff}
An \Ainf{}-category is an \Ainf{}-precategory together with a proof that its
operations satisfy the Stasheff identities.
\end{definition}

\begin{lemma}[Stasheff sum vanishes in an \Ainf{}-category]
\lean{AInfinityCategoryTheory.AInfinityCategory.stasheff_eq_zero}
\label{AInfinityCategoryTheory.AInfinityCategory.stasheff_eq_zero}
\leanok
\uses{AInfinityCategoryTheory.AInfinityCategory, AInfinityCategoryTheory.AInfinityPreCategory.stasheffSum}
For an \Ainf{}-category, every Stasheff sum computed from the structure maps is
equal to zero.
\end{lemma}
